\section*{Chapter \ref{ch:iv1:waiting-times-and-markov-chains} --- Waiting Times and Markov Chains}
\markboth{\protect\omTocMark{}{Answers: Waiting Times and Markov Chains}}{\protect\omTocMark{}{Answers: Waiting Times and Markov Chains}}

\begin{solution}{iq:iv1:waiting-times-and-markov-chains:1}
Two heads in a row takes 6 flips on average. A head followed by a tail takes 4.

Use first-step analysis\index{first-step analysis} on the states of \cref{fig:iv1:hh-ht}. For two heads, let $x$ be the expected number of flips from the start and $y$ the number after one head. A tail from either state sends you back to the start:
\[
x = 1 + \tfrac12 x + \tfrac12 y, \qquad y = 1 + \tfrac12 x,
\]
which gives $x = 6$ and $y = 4$. For a head then a tail, wait for the first head, 2 flips on average, and then for the first tail, 2 more. A further head after the first head loses nothing.

\begin{omfigure}
\begin{tikzpicture}[>=stealth,node distance=2.1cm,
  st/.style={circle,draw=omDef,thick,minimum size=11mm,inner sep=0pt,font=\small},
  lab/.style={font=\footnotesize,midway}]
\node[st] (a) {start};
\node[st,right of=a] (b) {H};
\node[st,right of=b,fill=omDef!15] (c) {HH};
\draw[->] (a) to[bend left=20] node[lab,above] {H} (b);
\draw[->] (b) to[bend left=20] node[lab,below] {T} (a);
\draw[->] (b) -- node[lab,above] {H} (c);
\draw[->] (a) to[out=135,in=225,looseness=5] node[lab,left] {T} (a);
\node[st,right of=c,xshift=1.2cm] (d) {start};
\node[st,right of=d] (e) {H};
\node[st,right of=e,fill=omDef!15] (f) {HT};
\draw[->] (d) -- node[lab,above] {H} (e);
\draw[->] (e) -- node[lab,above] {T} (f);
\draw[->] (d) to[out=135,in=225,looseness=5] node[lab,left] {T} (d);
\draw[->] (e) to[out=60,in=120,looseness=5] node[lab,above] {H} (e);
\end{tikzpicture}

\omcaption{Waiting for HH (left) and for HT (right). A wrong flip on the way to HH sends you back to the start, while a wrong flip on the way to HT keeps you where you are.}\label{fig:iv1:hh-ht}
\end{omfigure}

\iqwrong{``Each pattern has probability $1/4$ at any place, so both take 4 flips.'' Both do occur once per 4 flips in the long run, but occurrences of HH come in clumps, because HHH holds two of them. A clumpy pattern has longer gaps between clumps, so the first occurrence takes longer.}
\end{solution}

\begin{solution}{iq:iv1:waiting-times-and-markov-chains:2}
Two sixes in a row take 42 rolls on average. Three heads in a row take 14 flips.

Build up one at a time with first-step analysis\index{first-step analysis}. To get $k$ in a row, first get $k - 1$ in a row, which takes $E_{k-1}$ rolls, then roll once more. With probability $p$ you are done; otherwise you start again from nothing:
\[
E_k = E_{k-1} + 1 + (1 - p)\,E_k, \quad\text{so}\quad E_k = \frac{E_{k-1} + 1}{p}.
\]
For the die, $E_1 = 6$ and $E_2 = (6 + 1) \cdot 6 = 42$. For the coin, $E_1 = 2$, $E_2 = 6$ and $E_3 = 14$.

\iqwrong{36, the reciprocal of the chance of a pair of sixes at a given place. That is the long-run frequency, and it ignores the clumping that makes a first occurrence later.}
\iqfollow{Three sixes in a row?}{$(42 + 1) \cdot 6 = 258$. The general form is $6(6^k - 1)/5$.}
\end{solution}

\begin{solution}{iq:iv1:waiting-times-and-markov-chains:3}
Your probability of winning is $6/11 \approx 0.545$.

Use first-step analysis\index{first-step analysis} over a full round: you win now with probability $1/6$, or both of you miss, with probability $(5/6)^2$, and the game restarts as if new, so
\[
P = \tfrac16 + \bigl(\tfrac56\bigr)^2 P, \quad\text{which gives}\quad P = \frac{1/6}{11/36} = \frac{6}{11}.
\]
\iqfollow{You flip a coin with heads probability $1/6$ and your opponent one with $1/3$, you going first, and the first head wins. Your chance?}{$\frac{1/6}{1 - (5/6)(2/3)} = 3/8$. Going first does not make up for the weaker coin.}
\end{solution}

\enlargethispage{-3\baselineskip}%
\begin{solution}{iq:iv1:waiting-times-and-markov-chains:4}
A wins with probability $9/13 \approx 0.692$, and the game lasts $50/13 \approx 3.85$ more points on average.

Look at the game two points at a time, using first-step analysis\index{first-step analysis}. From a level score, A wins both with $p^2 = 0.36$, B wins both with $q^2 = 0.16$, and otherwise the players are level again, with probability $2pq = 0.48$. So
\[
W = p^2 + 2pq\,W \;\Rightarrow\; W = \frac{p^2}{p^2 + q^2} = \frac{0.36}{0.52} = \frac{9}{13},
\]
\[
L = 2 + 2pq\,L \;\Rightarrow\; L = \frac{2}{p^2 + q^2} = \frac{50}{13}.
\]
\iqwhy{The rule amplifies an edge: a 60\% point becomes a 69\% game. It is the simplest example of how repetition with a stopping rule turns a small edge into a large one, at the cost of a longer game.}
\end{solution}

\begin{solution}{iq:iv1:waiting-times-and-markov-chains:5}
The process spends a quarter of its time in B, a visit to B lasts $10/3$ steps on average, and the correlation is $0.6^n$.

For the stationary distribution\index{stationary distribution}, the flow from A to B must balance the flow back: $\pi_A \cdot 0.1 = \pi_B \cdot 0.3$, and so $\pi_A = 3/4$ and $\pi_B = 1/4$. A visit to B ends with probability $0.3$ at each step, so its length is geometric with mean $1/0.3 = 10/3$.

For the correlation, a two-state chain has eigenvalues 1 and $\lambda = 0.9 + 0.7 - 1 = 0.6$. The deviation from the stationary distribution shrinks by $\lambda$ at each step, so the lag-$n$ correlation is $\lambda^n = 0.6^n$.

\iqfollow{A bit passes through 10 copiers, and each flips it with probability $0.1$. What is the chance that it arrives correct?}{This is the same chain with $\lambda = 0.8$: $\tfrac12(1 + 0.8^{10}) \approx 0.554$.}
\iqwhy{This is the two-state regime model. Given the persistence of each state, it gives the share of time in the state, the length of a stay and the decay of any signal based on the state. Here the decay of $0.6$ per step means a half-life of little over one step, even though a stay in A lasts 10 steps.}
\end{solution}

\begin{solution}{iq:iv1:waiting-times-and-markov-chains:6}
Your probability of winning is $1/4$.

Think about when heads-heads-tails can win. If the first two flips are both heads, then you win for certain: the first tail that follows completes HHT, and THH needs a tail earlier. That has probability $1/4$. In every other case a tail appears before the first HH. Then the first HH to appear is preceded by a tail, and so THH is completed first.

\iqwrong{One half, by symmetry. The patterns are mirror images of each other, but a race reads flips forward only, so the symmetry between them fails.}
\iqfollow{And with a coin that shows heads with probability $2/3$?}{The argument is unchanged: HHT wins only by opening with HH, so its chance is $(2/3)^2 = 4/9$.}
\iqfollow{How long does the game last?}{$13/2$ flips on average, shorter than the 8 flips that either pattern needs alone.}
\end{solution}

\begin{solution}{iq:iv1:waiting-times-and-markov-chains:7}
HTHTH takes 42 flips on average, and HHHTT takes 32. The gap comes from how each pattern overlaps with itself.

Conway's formula\index{Conway's formula} says that the expected wait is the sum of $2^k$ over all $k$ such that the last $k$ letters equal the first $k$ letters. HTHTH overlaps itself at $k = 1, 3, 5$, which gives $2 + 8 + 32 = 42$. HHHTT overlaps only at $k = 5$, which gives 32.

To see why, use a casino argument\index{martingale casino argument}. Before each flip, a new gambler arrives with \$1 and bets it on the next flip being the pattern's first letter. A winner bets all of the \$2 on the second letter, and so on, until the gambler loses or completes the pattern. All bets are fair, so at the moment you stop, the casino's expected profit is zero: the $T$ dollars taken from gamblers equal in expectation what it pays out. At the end, the only gamblers still holding money are those who started $k$ flips ago, with the last $k$ flips being the first $k$ letters of the pattern. Each of those holds $2^k$. So $E[T] = \sum 2^k$ over those $k$.

\iqfollow{On a die, the pair 4 then 5 against the pair 4 then 4?}{$6^2 = 36$ rolls, against $36 + 6 = 42$.}
\iqfollow{A monkey types uniform capital letters. How long until ABRACADABRA?}{The word overlaps itself at lengths 1, 4 and 11, so $26^{11} + 26^4 + 26$ keystrokes.}
\end{solution}

\enlargethispage{-4\baselineskip}%
\begin{solution}{iq:iv1:waiting-times-and-markov-chains:8}
It takes 10 steps on average.

Use symmetry to lump the corners\index{state lumping} by distance from the target, as in \cref{fig:iv1:cube}. Let $a$, $b$, $c$ be the expected steps from distance 3, 2, 1. From distance 3, every edge leads to distance 2. From distance 2, two of three edges lead closer. From distance 1, one of three reaches the target:
\[
a = 1 + b, \qquad b = 1 + \tfrac13 a + \tfrac23 c, \qquad c = 1 + \tfrac23 b.
\]
These give $b = 9$, $c = 7$ and $a = 10$.

\begin{omfigure}
\begin{tikzpicture}[>=stealth,
  st/.style={circle,draw=omDef,thick,minimum size=9mm,inner sep=0pt,font=\small},
  lab/.style={font=\footnotesize,midway}]
\node[st] (a) at (0,0) {3};
\node[st] (b) at (3,0) {2};
\node[st] (c) at (6,0) {1};
\node[st,fill=omDef!15] (d) at (9,0) {0};
\draw[->] (a) to[bend left=20] node[lab,above] {$1$} (b);
\draw[->] (b) to[bend left=20] node[lab,below] {$\tfrac13$} (a);
\draw[->] (b) to[bend left=20] node[lab,above] {$\tfrac23$} (c);
\draw[->] (c) to[bend left=20] node[lab,below] {$\tfrac23$} (b);
\draw[->] (c) -- node[lab,above] {$\tfrac13$} (d);
\node[below=3mm of a,font=\footnotesize] {start};
\node[below=3mm of d,font=\footnotesize] {target};
\end{tikzpicture}

\omcaption{The ant's distance from the opposite corner. The eight corners of the cube reduce to four states, with one, three, three and one corners in them.}\label{fig:iv1:cube}
\end{omfigure}

\iqfollow{How long to return to the start?}{8 steps. The walk is uniform over 8 corners in the long run, and a state with long-run share $\pi$ recurs every $1/\pi$ steps on average.}
\iqfollow{How long to reach a neighbouring corner?}{7 steps. The commute time across an edge is $2 \cdot 12 \cdot \tfrac{7}{12} = 14$, and by symmetry half of it goes each way.}
\iqfollow{And on a four-dimensional cube?}{The same distance chain gives $64/3 \approx 21.33$.}
\end{solution}

\begin{solution}{iq:iv1:waiting-times-and-markov-chains:9}
On average it takes 6 rolls.

No first-step system is needed. The residue of the total modulo 6 is uniform after every roll, whatever came before, because the six faces reach each of the six residues exactly once from any starting residue. So each roll makes the total divisible by 6 with probability $1/6$, independently of the past, and the wait is geometric with mean 6.

\iqfollow{And divisible by 7?}{7 rolls. The first roll can never succeed, since it gives 1 to 6. After that, from any non-zero residue, exactly one face reaches a multiple of 7, and a failure leaves you at another non-zero residue. So it is one roll, plus a geometric wait with mean 6.}
\iqfollow{Why does this not work for 5?}{Five residues and six faces do not divide evenly, so the residue is not uniform and you need the first-step equations\index{first-step analysis}.}
\end{solution}

\enlargethispage{-3\baselineskip}%
\begin{solution}{iq:iv1:waiting-times-and-markov-chains:10}
To reach 10 or more takes about $3.32$ rolls. To reach 100 or more takes about $29.05$ rolls.

The exact method is a recursion by first-step analysis\index{first-step analysis}. Let $T(n)$ be the expected rolls to add at least $n$. One roll of $k$ leaves $n - k$ to go, and $T(n) = 0$ for $n \le 0$:
\[
T(n) = 1 + \frac16 \sum_{k=1}^{6} T(n - k).
\]
So $T(1) = 1$, $T(2) = 7/6$, $T(3) = 49/36$, and the table goes on to $T(10) \approx 3.32$.

For the check, use Wald's identity\index{Wald's identity}: the total when you stop averages $3.5\,E[T]$, and that total is $n$ plus an overshoot. For large $n$ the overshoot has mean
\[
\frac{E[X^2]}{2E[X]} - \frac12 = \frac{91/6}{7} - \frac12 = \frac53.
\]
So $E[T] \approx (n + 5/3)/3.5$, which gives $3.33$ for $n = 10$, and $29.05$ for $n = 100$.

\iqwrong{$n/3.5$, which is $2.86$ for $n = 10$. It forgets that you stop past the target, not on it.}
\end{solution}

\begin{solution}{iq:iv1:waiting-times-and-markov-chains:11}
The probability is $7/13 \approx 0.538$.

Each roll of the pair gives a 12 with probability $1/36$, a 7 with probability $1/6$, and something else with probability $29/36$. The only memory you need is whether the last sum was a 7. Let $a$ be the chance that 12 comes first when the last sum was not a 7, and $b$ the chance when it was. By first-step analysis\index{first-step analysis}, a 7 in the second state ends the game with a loss:
\[
a = \tfrac{1}{36} + \tfrac16 b + \tfrac{29}{36} a, \qquad b = \tfrac{1}{36} + \tfrac{29}{36} a.
\]
Substituting $b$ into the first equation and multiplying through by 216 gives $42a = 7 + 29a$, so $a = 7/13$.

\iqwrong{Comparing the chance of a 12 per roll, $1/36$, with the chance of two sevens, $1/36$, and answering $1/2$. Two consecutive sevens need two rolls that overlap, which makes them slower to arrive than a 12.}
\end{solution}

\begin{solution}{iq:iv1:waiting-times-and-markov-chains:12}
The expected wait is $6.25$ minutes, not 5.

This is the inspection paradox\index{inspection paradox}. A random moment tends to fall inside a long gap: three times as much clock time lies in 15-minute gaps as in 5-minute gaps. So you land in a 15-minute gap with probability $3/4$, and you then wait half of it on average:
\[
\tfrac34 \cdot 7.5 + \tfrac14 \cdot 2.5 = 6.25.
\]
The general formula is $E[G^2]/(2E[G])$, here $125/20$.

\iqfollow{Bus gaps that are exponential with mean 10 minutes?}{The wait is 10 minutes, the same as the whole mean gap. The gap you land in averages 20 minutes. For regular 10-minute buses the wait is 5.}
\iqfollow{In a long sequence of coin flips, how long is the run that contains flip number 1,000?}{3 on average, while a typical run has length 2. The same size-biasing applies: $E[L^2]/E[L] = 6/2$.}
\iqwhy{Any statistic sampled at random times sees long intervals too often. Measure the time between trades from a random clock and you overstate it, and measure the spread at random times and you over-weight quiet books. Decide whether you want an average per event or per unit of time.}
\end{solution}

\begin{solution}{iq:iv1:waiting-times-and-markov-chains:13}
The colony eventually dies out with probability $1/2$.

Use conditioning on the first step\index{branching processes}. Let $q$ be the probability of extinction. A family dies out only if each bacterium alive after the first minute starts a family that dies out, and those families are independent:
\[
q = \tfrac14 + \tfrac14 q + \tfrac12 q^2, \quad\text{or}\quad 2q^2 - 3q + 1 = 0.
\]
The roots are $q = 1/2$ and $q = 1$. The mean number of offspring is $\tfrac14 + 2 \cdot \tfrac12 = 5/4$, above 1, and in that case the extinction probability is the smaller root. If the mean is at most 1 (and a bacterium is not simply replaced each time), extinction is certain.

\iqwrong{``$q = 1$ solves the equation, so extinction is certain.'' It always solves it. The extinction probability is the smallest root in $[0, 1]$, which is the limit of the chance of dying out by each generation, starting from zero.}
\iqfollow{Starting from three bacteria?}{The families are independent, so $(1/2)^3 = 1/8$.}
\end{solution}

\begin{solution}{iq:iv1:waiting-times-and-markov-chains:14}
The most likely final total is exactly 1,000, with probability about $2/7 \approx 0.286$.

Far from the start, the chance that the running total ever lands on a given number settles at $u = 1/3.5 = 2/7$, by renewal theory\index{renewal argument}: each roll advances the total by $3.5$ on average, so in the long run one number in $3.5$ gets visited.

To finish at $1000 + k$, for $k = 0, \dots, 5$, you must pass through one of the numbers in $1000 + k - 6, \dots, 999$, and then roll exactly the right face. There are $6 - k$ such numbers, each visited with chance $2/7$, and each needs one face, with chance $1/6$. So
\[
P(\text{final total} = 1000 + k) = (6 - k) \cdot \frac27 \cdot \frac16 = \frac{6 - k}{21}.
\]
At $k = 0$ this is $6/21 = 2/7$, the largest. The chances are $6/21, 5/21, \dots, 1/21$, falling to the last at 1,005. Their mean overshoot is $5/3$.

\iqwrong{``Uniform over 1,000 to 1,005, so $1/6$.'' That would need each of the final totals to be reached equally often, but 1,000 can be reached from six numbers, and 1,005 from only one.}
\iqfollow{What is the chance that the roll crossing 1,000 is a 2?}{By the same count, $2 \cdot \tfrac27 \cdot \tfrac16 = 2/21$. Large rolls are more likely to be the crossing one: the chance for a roll of $j$ is $j/21$.}
\end{solution}

\begin{solution}{iq:iv1:waiting-times-and-markov-chains:15}
Against HHH, name THH and you win with probability $7/8$. Against HTH, name HHT and win with $2/3$. Whatever your opponent names, there is a response that gives at least $2/3$.

This is Penney's game\index{Penney's game}. The recipe is to take your opponent's first two letters as your last two, and to begin with the opposite of their second letter. HHH wins only if the first three flips are all heads, since any tail before that sets up THH first. That gives THH the chance $7/8$.

In general, use Conway's leading numbers\index{Conway's formula}: $L(X, Y)$ is the sum of $2^{k-1}$ over all $k$ for which the last $k$ letters of $X$ equal the first $k$ of $Y$. The odds that $B$ beats $A$ are then
\[
\bigl(L(A, A) - L(A, B)\bigr) : \bigl(L(B, B) - L(B, A)\bigr).
\]
For $A = $ HTH and $B = $ HHT, $L(A, A) = 5$, $L(A, B) = 1$, $L(B, B) = 4$ and $L(B, A) = 2$, so the odds are $4:2$, which is $2/3$.

The complete table of best responses is: HHH to THH ($7/8$), HHT to THH ($3/4$), HTH to HHT ($2/3$), HTT to HHT ($2/3$), and the same with heads and tails swapped.

\iqwrong{``Pick a pattern that appears as often, so it is fair.'' All eight appear equally often. The relation is not even transitive: THH beats HHT, which beats HTT, which beats TTH, which beats THH.}
\end{solution}

\begin{solution}{iq:iv1:waiting-times-and-markov-chains:16}
The probability is $1/11$. The whole walk takes $66$ steps on average.

Every number other than the start is equally likely to be visited last. For the 6 to be last, the walk must first reach one of its two neighbours, 5 or 7; it must, because it visits every number. Say it reaches 5 first. Then the 6 is last exactly if the walk goes from 5 all the way round to 7 before stepping on 6. On the line of positions from 6 round to 7, the walk stands one step from one end. A gambler's ruin\index{gambler's ruin} on a fair walk with 11 steps between the ends gives the chance of reaching the far end as $1/11$. Nothing in this depends on which number you chose, so all 11 have the same chance.

For the steps, at any time the visited numbers form an arc. With an arc of $j$ numbers you stand at one end of it, and the next new number takes $j$ steps on average: it is a fair walk on $j + 2$ positions that starts next to one end, and such a walk takes $1 \cdot j$ steps. Summing from $j = 1$ to 11 gives $66$, or $n(n-1)/2$ in general.

\iqwrong{Assuming the 6, being farthest away, is most likely to be last. It takes longest to reach on average, at $6 \cdot 6 = 36$ steps, but that is a different question.}
\end{solution}

\begin{solution}{iq:iv1:waiting-times-and-markov-chains:17}
It takes 168 moves on average.

No calculation over the board is needed. For a random walk on a graph, which picks a uniformly random edge at each step, the stationary distribution\index{stationary distribution} is proportional to the degree, the number of edges at a vertex:
\[
\pi(v) = \frac{\deg(v)}{\sum_u \deg(u)}.
\]
You can verify this in one line: the flow into $v$ is $\sum_{u \sim v} \pi(u)/\deg(u)$, which has $\deg(v)$ equal terms, each $1/\sum_u \deg(u)$. The expected return time to a state is $1/\pi(v)$, by Kac's lemma\index{Kac's lemma}.

Now count the knight's moves. A knight move has 2 orientations and 4 sign choices, and each of the 8 kinds has $7 \cdot 6$ placements on the board, so $\sum_u \deg(u) = 8 \cdot 42 = 336$. A corner has degree 2, so the return takes $336/2 = 168$ moves.

\iqwrong{Trying to solve 63 linear equations on the board. An interviewer asking this wants the degree argument.}
\iqfollow{From a central square?}{Degree 8, so $336/8 = 42$ moves.}
\iqfollow{A king starting in a corner?}{The king's degrees add up to 420, and a corner has degree 3, so 140 moves.}
\end{solution}

\begin{solution}{iq:iv1:waiting-times-and-markov-chains:18}
It takes 18 steps on average.

Reflection across the line from start to target lumps the positions\index{state lumping} into classes. First-step analysis\index{first-step analysis} on them gives the values in \cref{fig:iv1:grid}:
\[
s = 1 + a, \quad a = 1 + \tfrac13(s + m + c), \quad m = 1 + \tfrac14(2a + 2b), \quad c = 1 + \tfrac12(a + b), \quad b = 1 + \tfrac13(m + c),
\]
where $s$, $a$, $m$, $c$, $b$ are the start, its neighbours, the centre, the two other corners and the target's neighbours. The equations for $m$ and $c$ have the same right-hand side, so $m = c$. The solution is $s = 18$, $a = 17$, $m = c = 15$ and $b = 11$.

\begin{omfigure}
\begin{tikzpicture}[x=1.7cm,y=1.7cm,
  st/.style={circle,draw=omDef,thick,minimum size=8mm,inner sep=0pt,font=\small,fill=white}]
\foreach \i in {0,1,2}{\draw[omDef] (\i,0) -- (\i,2); \draw[omDef] (0,\i) -- (2,\i);}
\node[st] at (0,0) {18};
\node[st] at (1,0) {17};
\node[st] at (2,0) {15};
\node[st] at (0,1) {17};
\node[st] at (1,1) {15};
\node[st] at (2,1) {11};
\node[st] at (0,2) {15};
\node[st] at (1,2) {11};
\node[st,fill=omDef!15] at (2,2) {0};
\node[below=5mm,font=\footnotesize] at (0,0) {start};
\node[above=5mm,font=\footnotesize] at (2,2) {target};
\end{tikzpicture}

\omcaption{Expected steps to reach the top-right corner from each position of the grid. Each value is one more than the average over its neighbours.}\label{fig:iv1:grid}
\end{omfigure}

A neater route is the commute time identity\index{commute time identity}: treat each connection as a 1-ohm resistor. The expected round trip between two vertices is then $2 \cdot (\text{number of edges}) \cdot R$, with $R$ the effective resistance. Here $R = 3/2$, so a round trip takes $2 \cdot 12 \cdot 3/2 = 36$ steps. By symmetry each way takes half of that.

\iqfollow{Is the time always the same each way?}{No. From centre to corner takes 15 steps, and from corner to centre only 6. Their sum, 21, is $2 \cdot 12 \cdot 7/8$, as the identity says.}
\end{solution}

\begin{solution}{iq:iv1:waiting-times-and-markov-chains:19}
The variance is 22, so the standard deviation is about $4.7$, close to the mean of 6. For a head then a tail, the mean is 4 and the variance is only 4.

Use first-step analysis\index{first-step analysis} on the second moment\index{second moment by first step}. With means $m_0 = 6$ from the start and $m_1 = 4$ after a head, let $a = E[T_0^2]$ and $b = E[T_1^2]$. One flip adds one to the remaining time, and $E[(1 + T)^2] = 1 + 2E[T] + E[T^2]$. So
\[
a = \tfrac12(1 + 2 \cdot 6 + a) + \tfrac12(1 + 2 \cdot 4 + b) = 11 + \tfrac12 a + \tfrac12 b,
\]
\[
b = \tfrac12 \cdot 1 + \tfrac12(1 + 2 \cdot 6 + a) = 7 + \tfrac12 a.
\]
This gives $a = 58$ and $b = 36$, so the variance is $58 - 6^2 = 22$.

The wait for a head then a tail is the sum of two independent geometric waits, each with mean 2 and variance 2, so it has variance 4.

\iqfollow{And for three heads in a row?}{The same method gives a mean of 14 and a variance of 142.}
\iqwhy{Two waits with similar means can have very different risk. A wait whose failures send it back to the beginning has a standard deviation close to its mean, so plan for a long tail and do not quote the mean alone.}
\end{solution}

\begin{solution}{iq:iv1:waiting-times-and-markov-chains:20}
It takes $52\,H_{52} \approx 236$ moves on average, where $H_n = \sum_{k \le n} 1/k$.

Follow the card that started at the bottom. When $i$ cards are below it, the top card lands below it with probability $(i + 1)/52$, because there are $i + 1$ positions below it. So the wait for that card to move up one place has mean $52/(i + 1)$. It reaches the top when $i = 51$, and one more move inserts it:
\[
\sum_{i=0}^{50} \frac{52}{i+1} + 1 = 52\,H_{51} + 1 = 52\,H_{52} \approx 236.
\]
This is the coupon collector\index{coupon collector} in another guise.

The deck is perfectly shuffled at that moment by induction\index{strong uniform time}. The cards below the marked card are always in uniformly random order: a new card goes to a uniformly random one of the $i + 1$ places among the $i$ cards already there. When the marked card reaches the top, the 51 cards beneath it are in uniform order. The final insertion places it uniformly among the 52 positions, so the entire deck is uniform.

\iqfollow{Is 236 the mixing time?}{It is an upper bound in expectation: the deck may already be close to uniform earlier. It tells you what order of work to expect: $n \ln n$ moves.}
\end{solution}

\begin{solution}{iq:iv1:waiting-times-and-markov-chains:21}
Reaching all heads takes $74752/63 \approx 1186.5$ steps on average. Coming back to all tails takes exactly 1,024.

The return time needs almost no algebra. The state is a corner of a 10-dimensional cube, and the walk has a uniform stationary distribution\index{stationary distribution} over the $2^{10}$ corners, because each corner has the same degree. So by Kac's lemma\index{Kac's lemma} the mean return time is $2^{10} = 1024$.

For reaching all heads, track only the number of tails. Let $e_k$ be the expected steps to go from $k$ tails to $k - 1$. With probability $k/10$ you flip a tail and are done. Otherwise you create a tail, and you must come down from $k + 1$ and then from $k$ again:
\[
e_k = 1 + \frac{10 - k}{10}(e_{k+1} + e_k) \quad\Rightarrow\quad e_k = \frac{10}{k} + \frac{10 - k}{k}\,e_{k+1},
\]
with $e_{10} = 1$. Working downwards and summing $e_{10} + \dots + e_1$ gives $74752/63$. The same recursion with 3 coins gives 10, the ant on the cube.

\iqwrong{Assuming the far state needs far longer. Both times are of order $2^{10}$, because the bulk of the time is spent wandering around the middle, near five heads, before a rare run to either extreme.}
\end{solution}

\begin{solution}{iq:iv1:waiting-times-and-markov-chains:22}
Name THTH. You win with probability $9/14 \approx 0.643$, and the game lasts $90/7 \approx 12.86$ flips on average.

Use Conway's leading numbers\index{Conway's formula} with two casinos\index{martingale casino argument}. As with one pattern, a gambler who started $k$ flips ago holds $2^k$ if the last $k$ flips match the first $k$ letters of the pattern that gambler is betting on. Run one casino of gamblers betting on $A$ and another betting on $B$. Let $L(X, Y) = \sum 2^{k-1}$ over the $k$ for which the last $k$ letters of $X$ are the first $k$ of $Y$, and let $p$ be the chance that $A$ wins. Each casino breaks even in expectation, so
\[
E[T] = 2p\,L(A,A) + 2(1-p)\,L(B,A) = 2p\,L(A,B) + 2(1-p)\,L(B,B).
\]
Equating the two right-hand sides gives the odds of $B$ against $A$, $\;L(A,A) - L(A,B) : L(B,B) - L(B,A)$.

The best responses start with one letter and then copy the first three of $A$, here $x$HTH, so that $B$ needs one flip less whenever $A$ is under way. For $A = $ HTHH and $B = $ THTH, $L(A,A) = 9$, $L(A,B) = 0$, $L(B,B) = 10$ and $L(B,A) = 5$. The odds are $9:5$, so $9/14$. Then $E[T] = 2 \cdot \tfrac{5}{14} \cdot 9 + 2 \cdot \tfrac{9}{14} \cdot 5 = 90/7$. The second casino gives the same: $2 \cdot \tfrac{9}{14} \cdot 10 = 90/7$. The other candidate, HHTH, scores only $(9-3):(9-5)$, or $3/5$.

\iqwrong{Using the three-flip recipe without checking: begin with the opposite of the second letter of $A$, then copy the first three letters of $A$. It gives HHTH, and that is not the best. For four flips the first letter must be chosen by computing both.}
\end{solution}

\begin{solution}{iq:iv1:waiting-times-and-markov-chains:23}
The expected number of rolls is $416/5 = 83.2$.

Find the right states\index{state lumping}. Let the state $j$ be the number of most recent rolls that are all different, so you start at state 1 after one roll and stop at state 6. From state $j$, a new face has probability $(6 - j)/6$ and moves you to $j + 1$. Otherwise the roll repeats the face that appeared $i$ rolls ago, for one of $i = 1, \dots, j$, each with probability $1/6$, and the new distinct tail has length exactly $i$, as in \cref{fig:iv1:distinct}.

Let $e_j$ be the expected time to first reach $j + 1$ from $j$. A drop to state $i$ costs the climb back, $e_i + \dots + e_j$. Writing out the first step and collecting terms,
\[
e_j = 1 + \frac16 \sum_{i=1}^{j} \sum_{m=i}^{j} e_m = 1 + \frac16 \sum_{m=1}^{j} m\,e_m \quad\Rightarrow\quad (6 - j)\,e_j = 6 + \sum_{m < j} m\,e_m.
\]
This solves from the bottom up: $e_1 = \tfrac65$, $e_2 = \tfrac95$, $e_3 = \tfrac{18}{5}$, $e_4 = \tfrac{54}{5}$, $e_5 = \tfrac{324}{5}$. Adding the first roll, the total is $1 + 411/5 = 416/5$.

\begin{omfigure}
\begin{tikzpicture}[>=stealth,x=1.8cm,
  st/.style={circle,draw=omDef,thick,minimum size=8mm,inner sep=0pt,font=\small},
  lab/.style={font=\footnotesize}]
\foreach \j in {1,...,5}{\node[st] (s\j) at (\j,0) {\j};}
\node[st,fill=omDef!15] (s6) at (6,0) {6};
\draw[->] (s1) -- (s2);
\draw[->] (s2) -- (s3);
\draw[->] (s3) -- (s4);
\draw[->] (s4) -- node[above,lab] {$\tfrac26$} (s5);
\draw[->] (s5) -- (s6);
\draw[->,omThm] (s4) to[bend left=45] node[below,lab,pos=0.5] {$\tfrac16$} (s1);
\draw[->,omThm] (s4) to[bend left=50] node[above,lab,pos=0.5] {$\tfrac16$} (s2);
\draw[->,omThm] (s4) to[bend left=60] node[above,lab,pos=0.5] {$\tfrac16$} (s3);
\draw[->,omThm] (s4) to[out=60,in=120,looseness=5] node[above,lab] {$\tfrac16$} (s4);
\end{tikzpicture}

\omcaption{The chain on the number of distinct recent rolls, showing only the moves from state 4. A repeat of the face seen $i$ rolls ago drops you to state $i$. Every state behaves the same way.}\label{fig:iv1:distinct}
\end{omfigure}

\iqwrong{$6^6/6! = 64.8$, the reciprocal of the chance that a given window of six is all different. That is the long-run frequency of such windows, and the windows overlap and clump, so the first one comes later.}
\end{solution}

\begin{solution}{iq:iv1:waiting-times-and-markov-chains:24}
The family dies out with probability $1/2$, and given that, the expected total is 3 cells.

The extinction probability solves $q = \tfrac13 + \tfrac23 q^2$, by conditioning on the first step\index{branching processes}. The roots are $1/2$ and 1, and with mean offspring $4/3$ above 1 the answer is the smaller root.

Unconditionally, the expected total is infinite, because the family grows without limit half the time. The idea is to condition on extinction\index{conditioning on extinction}. A cell that has $k$ children has a family that dies out with probability $q^k$, independently for each child, so the conditional law of the number of children is $P(k \text{ children} \mid \text{extinct}) = p_k\,q^k / q$:
\[
P(0 \mid \text{extinct}) = \frac{1/3}{1/2} = \frac23, \qquad P(2 \mid \text{extinct}) = \frac{(2/3)(1/4)}{1/2} = \frac13.
\]
Each child's family, given that it dies out, is again a copy of the same conditional process. So conditional on extinction you have a branching process with mean offspring $\tfrac23$. Generation $n$ has expected size $(2/3)^n$, and the total is
\[
\sum_{n \ge 0} \Bigl(\frac23\Bigr)^n = 3.
\]
\iqwrong{Saying that the total is infinite, or computing the mean with the original probabilities. Conditioning changes the offspring law, and this is the step that makes the question hard.}
\iqfollow{Given extinction, what is the chance the first cell dies at once?}{$2/3$, the first probability above, up from $1/3$.}
\end{solution}

\begin{solution}{iq:iv1:waiting-times-and-markov-chains:25}
Generation 10 is empty with probability $10/11$. If it is not empty, its expected size is 11. The expected generation of extinction is infinite, although extinction is certain.

Use the probability generating function\index{probability generating function}. The offspring law has $f(s) = \sum_k s^k/2^{k+1} = 1/(2 - s)$, and generation $n$ has the $n$-fold composition $f_n$ of $f$. The idea that removes the algebra is that this map is linear after a change of variable\index{linear-fractional iteration}:
\[
\frac{1}{1 - f(s)} = \frac{2 - s}{1 - s} = 1 + \frac{1}{1 - s}.
\]
So each iteration adds 1, and $\dfrac{1}{1 - f_n(s)} = n + \dfrac{1}{1 - s}$. At $s = 0$ this gives $1/(1 - f_n(0)) = n + 1$, and so $P(\text{generation } n \text{ is empty}) = f_n(0) = n/(n+1)$. For $n = 10$ that is $10/11$.

The mean offspring is 1, so $E[X_n] = 1$ for every $n$. Given survival, $E[X_n \mid X_n > 0] = 1/P(X_n > 0) = n + 1$, which is 11.

The family outlasts generation $n$ with probability $1/(n + 1)$. Summing these tails for the expected extinction time gives the harmonic series, which diverges.

\iqwrong{``Mean offspring of one means a stable population.'' A critical family dies out with certainty. The mean stays at 1 only because the rare survivors grow: at generation $n$ only one family in $n + 1$ survives, with size $n + 1$ on average.}
\iqwhy{Anything that reproduces at exactly its replacement rate behaves like this: most lines die early, and the few survivors look large. Reading only the survivors overstates growth by a factor that rises with time.}
\end{solution}
